% BHU PYQ -- Professional Exam Template
\documentclass[11pt,a4paper]{article}

\usepackage[a4paper,top=2.5cm,bottom=2.5cm,left=2.8cm,right=2.8cm,headheight=14pt]{geometry}
\usepackage{amsmath,amssymb}
\usepackage[T1]{fontenc}
\usepackage[utf8]{inputenc}
\usepackage{lmodern}
\usepackage{microtype}
\usepackage{fancyhdr}
\usepackage{enumitem}
\usepackage{hyperref}
\usepackage{lastpage}
\usepackage{parskip}
\usepackage{tikz}
\usetikzlibrary{arrows.meta,shapes,positioning}

\hypersetup{colorlinks=true,urlcolor=black,linkcolor=black,
  pdftitle={B.Sc. (Hons.) Zoology Sem-I ZOB-101 (Old Course 2009-10) -- BHU PYQ}}

\pagestyle{fancy}
\fancyhf{}
\renewcommand{\headrulewidth}{0.4pt}
\renewcommand{\footrulewidth}{0.4pt}
\fancyhead[L]{\small   \textit{Banaras Hindu University}}
\fancyhead[R]{\small   \textit{Zoology Paper: ZOB-101 (Old 2009)}}
\fancyfoot[C]{\small Page  \thepage\ of \pageref{LastPage}}


\newlist{parts}{enumerate}{3}
\setlist[parts,1]{label=(\alph*),leftmargin=*,itemsep=4pt,topsep=3pt}
\setlist[parts,2]{label=(\roman*),leftmargin=*,itemsep=2pt,topsep=2pt}
\setlist[parts,3]{label=(\arabic*),leftmargin=*,itemsep=2pt,topsep=2pt}

\newcommand{\pts}[1]{\hfill   \textit{[#1]}}


\begin{document}


\begin{center}
  {\large\bfseries BANARAS HINDU UNIVERSITY}\\[2pt]
  {\normalsize B.Sc. (Hons.) Semester I Examination, 2016-17}\\[6pt]
  \rule{0.8\linewidth}{0.5pt}\\[6pt]
  {\large\bfseries Zoology}\\[4pt]
  {\large\bfseries Paper: ZOB-101 --- Systematics \& Animal Diversity and Animal Form \& Function}\\[2pt]
  {\normalsize   \textbf{(Old Course 2009-10 Syllabus)}}\\[4pt]
  {\normalsize Time: Three hours \hfill Full Marks: 70}
\end{center}

\rule{\linewidth}{0.4pt}

\smallskip



\noindent   \textit{Note: Answer three questions from Section-A and three questions from Section-B. Question No. 1 in both sections is compulsory. Use a separate answer book for each Section. The figures in the right-hand margin indicate marks.}

\bigskip

\begin{center}
   \textbf{SECTION -- A} \\
   \textbf{(Systematics \& Animal Diversity) (Marks: 35)}
\end{center}

\smallskip



\noindent   \textit{Note: Answer three questions including Question No. 1 which is compulsory.}

\medskip


\noindent   \textbf{Question 1.} Answer the following parts:

\begin{parts}
  \item Differentiate between the following:\pts{$2  \times 2 = 4$}
  
\begin{parts}
    \item   \textit{Petromyzon} and   \textit{Myxine}
    \item Prototheria and Eutheria.
  \end{parts}

  \item Define the following:\pts{$1  \times 3 = 3$}
  
\begin{parts}
    \item Binomial Nomenclature
    \item Species
    \item Homology.
  \end{parts}

  \item Write whether the following statements are True or False, giving reasons in two to three sentences:\pts{$2  \times 2 = 4$}
  
\begin{parts}
    \item Honey bees belong to class Crustacea of phylum Arthropoda.
    \item Water vascular system is found in   \textit{Asterias}.
  \end{parts}
\end{parts}

\medskip


\noindent   \textbf{Question 2.} Answer the following parts:\pts{$6 + 6 = 12$}

\begin{parts}
  \item Discuss the distinguishing features of lizards and snakes with at least three examples of each group.
  \item What is species concept? Discuss typological and biological species concept.
\end{parts}

\medskip


\noindent   \textbf{Question 3.} Answer the following parts:\pts{$6 + 6 = 12$}

\begin{parts}
  \item Enumerate salient features of the animals belonging to Phylum Mollusca. Classify this phylum up to its different classes with their examples.
  \item Describe flight adaptations in birds.
\end{parts}

\medskip


\noindent   \textbf{Question 4.} Write notes on the following:\pts{$2  \times 6 = 12$}

\begin{parts}
  \item Hemichordates
  \item Dipnoi.
\end{parts}

\pagebreak

\begin{center}
   \textbf{SECTION -- B} \\
   \textbf{(Animal Form \& Function) (Marks: 35)}
\end{center}

\smallskip



\noindent   \textit{Note: Answer three questions including Question No. 1 which is compulsory.}

\medskip


\noindent   \textbf{Question 1.} Answer the following parts:

\begin{parts}
  \item Define the following:\pts{$1  \times 3 = 3$}
  
\begin{parts}
    \item Food vacuole
    \item Malpighian tubules
    \item Microglia.
  \end{parts}

  \item State whether the following statements are True or False giving reasons:\pts{$2  \times 4 = 8$}
  
\begin{parts}
    \item Closed circulatory system is present in Earthworm.
    \item Flame cells are present in vertebrates for excretion.
    \item Bidder's canal is present in aves.
    \item Axon terminals are involved in the formation of synapse.
  \end{parts}
\end{parts}

\medskip


\noindent   \textbf{Question 2.} Answer the following parts:\pts{$6 + 6 = 12$}

\begin{parts}
  \item With the help of a suitable diagram describe the function of mammalian kidney.
  \item Explain the circulatory system of a cockroach.
\end{parts}

\medskip


\noindent   \textbf{Question 3.} Answer the following parts:\pts{$6 + 6 = 12$}

\begin{parts}
  \item Draw a labeled diagram of ommatidium and write its function.
  \item Discuss different types of asexual reproduction in animals.
\end{parts}

\medskip


\noindent   \textbf{Question 4.} Write short notes on the following:\pts{$3  \times 4 = 12$}

\begin{parts}
  \item Structure and function of gills
  \item Parthenogenesis
  \item Phonoreception in fish.
\end{parts}

\vfill
\rule{\linewidth}{0.4pt}

\begin{center}\small
\href{https://bhu-exam.vercel.app/}{Click Here}\\[4pt]\href{https://me-aryan.vercel.app/}{Click Here}\\[4pt]\href{https://quashan-me.vercel.app/}{Click Here}
\end{center}

\end{document}
